%%%%%%%%%%%%%%%%%%%%%%%%%%%%%%%%%%%%%%%%%%%%%%%%%%%%%%%%%%%%%%%%%%%%%%%%%
%                                                                       %
%                          TEMPLATE                                     %
%                                                                       %
%%%%%%%%%%%%%%%%%%%%%%%%%%%%%%%%%%%%%%%%%%%%%%%%%%%%%%%%%%%%%%%%%%%%%%%%%


\documentclass[12pt]{article}

%---------------------------------------------------------------------
% AMS packages 
%---------------------------------------------------------------------
 
% Packiet loaded automatically by amsart:
%  1. amsmath, 2. amsthm, 3. amsfonts.
 
%\usepackage{amssymb}
%\usepackage{amsmath}
 
% *** 'amsfonts': 
% boldface of the symbol of the real line: 'R'
\usepackage{amsfonts}
 
% *** 'amsthm': 
% 1. makes easy to modify the macro: \newtheorem{}{} since contains
% \theoremstyle{}
% 2. contains macro: \begin{proof} ... \end{proof}
\usepackage{amsthm} 


%%%%%%%%%%%%%%% Polish letter packages %%%%%%%%%%%%%%%%%%%%%%%%%%%%%%%%
%\usepackage[polish]{babel}
%\usepackage[latin2]{inputenc}
%\usepackage{t1enc}
 
%---------------------------------------------------------------------
% Theorems 
%---------------------------------------------------------------------
 
% Theorem style 'plain' are for: Theorem, Lemma, Corollary, 
% Proposition, Conjecture, Criterion, Algorithm  
\theoremstyle{plain} 
\newtheorem{theorem}{Theorem}[section] 
\newtheorem{lemma}[theorem]{Lemma} 
\newtheorem{question}[theorem]{Question}
\newtheorem{corollary}[theorem]{Corollary} 
\newtheorem{conclusion}[theorem]{Corollary} 
\newtheorem{claim}[theorem]{Claim} 
\newtheorem{fact}[theorem]{Fact} 
\newtheorem{proposition}[theorem]{Proposition} 
\newtheorem{axiom}{Axiom} 
 
% Theorem style 'definition' are for: Definition, Condition, Problem, 
% Example 
 
\theoremstyle{definition} 
\newtheorem{definition}[theorem]{Definition} 
\newtheorem{example}[theorem]{Example} 
\newtheorem{exercise}{Exercise} 
\newtheorem*{solution}{Solution} 
 
% Theorem style 'remark' are for: Remark, Note, Notation, Claim, 
% Summary, Acknowledgement, Case, Conclusion 
 
\theoremstyle{remark} 
\newtheorem*{remark}{Remark} 
\newtheorem*{notation}{Notation} 
\newtheorem*{acknowledgement}{Acknowledgement} 

%%%%%%%%%%%%%%%%%  The previous one %%%%%%%%%%%%%%%%%%%%%%%%%%%%%%%%%%%%%%
%\newtheorem{theorem}{Theorem}[section]
%\newtheorem{lemma}[theorem]{Lemma}
%\newtheorem{observation}[theorem]{Observation}
%\newtheorem{conclusion}[theorem]{Conclusion}
%\newtheorem{question}[theorem]{Question}
%\newtheorem{problem}[theorem]{Problem}
%\newtheorem{corollary}[theorem]{Corollary}
%\newtheorem{definition}[theorem]{Definiton}
%\newtheorem{claim}[theorem]{Claim}
%\newtheorem{fact}[theorem]{Fact}
%\newtheorem{conjecture}[theorem]{Conjecture}
%\newtheorem{example}[theorem]{Example}
%\newtheorem{proposition}[theorem]{Proposition}
%\newtheorem{remark}[theorem]{Remark}
%%%%%%%%%%%%%%%%%%%%%%%%%%%%%%%%%%%%%%%%%%%%%%%%%%%%%%%%%%%%%%%%%%%%%


%A
\newcommand{\afc}{AFC}
\newcommand{\afcbar}{\overline{AFC}}
\newcommand{\arr}{\rightarrow}
\newcommand{\Arr}{\Rightarrow}
\newcommand{\bp}{\begin{proof}}
\newcommand{\ep}{\end{proof}}

%B
\newcommand{\baire}{\omega^{\omega}}
\newcommand{\blacksquare}{\begin{flushright} \rule{2mm}{2mm} \end{flushright}}
\newcommand{\Bor}{\mbox{${\cal B}or$}}
%Previous seems to be much finer than next.
%\newcommand{\Bor}{{\it Bor}}
\newcommand{\borelucrz}{Borel-UCR_0}

%C
\newcommand{\ca}{2^{\omega}}
\newcommand{\cantor}{\ca}
\newcommand{\Card}[1]{\Vert #1 \Vert}

%D
\newcommand{\dom}{{\rm dom}}
\newcommand{\dummy}{{\tt Blah blah blah}}

%E
\newcommand{\Even}{\hbox{\rm \tiny Even}}

%a
\newcommand{\finsub}{[\omega]^{<\omega}}
\newcommand{\forces}{\mathrel{\|}\joinrel\mathrel{-}}

%G
\newcommand{\Graph}{\hbox{\it Graph}}

%H
\newcommand{\homeomorphic}{\approx}

%I
\newcommand{\incr}{\omega^{\uparrow \omega }}
\newcommand{\infsub}{[\omega]^{\omega}}

%L
\newcommand{\la}{\langle}

%M
\newcommand{\meager}{{\cal MGR}}
\newcommand{\minideal}{${\cal F}_{\hbox{\rm \scriptsize min}}(\neg
  D)\;$}

%N
\newcommand{\neglig}{{\cal N}}
\newcommand{\nnatural}{\mathbb{N}}

%O
\newcommand{\Odd}{\hbox{\rm \tiny Odd}}

%P
\newcommand{\Part}{{\it Part}}
\newcommand{\Perf}{{\it Perf}}
%%%\newcommand{\proof}{\flushleft{ \sc Proof. } \\ }
\newcommand{\Proof}[1]{\bigbreak\noindent{\bf Proof #1}\enspace}

%Q
%%%\newcommand{\qed}{{\hfill\vrule height6pt width6pt depth1pt\medskip}}
%\newcommand{\qed}{\sharp}
\newcommand{\QED}{\hspace{0.1in} \Box \vspace{0.1in}}

%R
\newcommand{\ra}{\rangle}
\newcommand{\ran}{{\rm ran}}
\newcommand{\rational}{\mathbb{Q}}
\newcommand{\real}{\mathbb{R}}
\newcommand{\restriction}{{\hbox{$\,|\grave{}\,$}}}

%S
\newcommand{\seq}{\subseteq}
%%%\newcommand{\square}{\hbox{\ \ \ \ \ \vrule\vbox{\hrule\phantom{o}\hrule}\vrule}}
% a small restriction:
\newcommand{\srestriction}{{\hbox{${\scriptstyle\,|\grave{}\,}$}}}

%U
\newcommand{\up}{\uparrow}
\newcommand{\ucrz}{UCR_0}

%%%%%%%%%%%%%%%%%%%%%% Calligraphic font commands %%%%%%%%%%%%%%%%%%%%%%%%%%
\newcommand{\cA}{{\cal A}}
\newcommand{\cB}{{\cal B}}
\newcommand{\cC}{{\cal C}}
\newcommand{\cD}{{\cal D}}
\newcommand{\cE}{{\cal E}}
\newcommand{\cF}{{\cal F}}
\newcommand{\cG}{{\cal G}}
\newcommand{\cH}{{\cal H}}
\newcommand{\cI}{{\cal I}}
\newcommand{\cJ}{{\cal J}}
\newcommand{\cK}{{\cal K}}
\newcommand{\cL}{{\cal L}}
\newcommand{\cM}{{\cal M}}
\newcommand{\cN}{{\cal N}}
\newcommand{\cO}{{\cal O}}
\newcommand{\cP}{{\cal P}}
\newcommand{\cQ}{{\cal Q}}
\newcommand{\cR}{{\cal R}}
\newcommand{\cS}{{\cal S}}
\newcommand{\cT}{{\cal T}}
\newcommand{\cU}{{\cal U}}
\newcommand{\cV}{{\cal V}}
\newcommand{\cW}{{\cal W}}
\newcommand{\cX}{{\cal X}}
\newcommand{\cY}{{\cal Y}}
\newcommand{\cZ}{{\cal Z}}

%%%%%%%%%%%%%%%%%%%%%% Some Greek fonts %%%%%%%%%%%%%%%%%%%%%%%%%%%%%%%%%%%%%%%
\newcommand{\oo}{\omega}
\newcommand{\bb}{\beta}
\newcommand{\dd}{\delta}
\newcommand{\ee}{\varepsilon}
\newcommand{\kk}{\kappa}
%%%\newcommand{\th}{\theta}

\renewcommand\abstractname{Summary}

%%%To print the date and time on each page
%%% comment out if not needed (next 14 lines)
%\makeatletter
%{\newcount\@hour}
%{\newcount\@minute}
%\def\timenow{\@hour=\time \divide\@hour by 60
%\number\@hour:
%  \multiply\@hour by 60 \@minute=\time
%  \global\advance\@minute by -\@hour
%  \ifnum\@minute<10 0\number\@minute\else
%  \number\@minute\fi}
%\def\ctimenow{\hfil{\tt \jobname.tex, \today~Time: \timenow }\hfil}
%      \let\@oddfoot\ctimenow\let\@evenfoot\ctimenow
%\makeatother

\pagestyle{myheadings}
%\markboth{{\bf Andrzej Nowik}}{\bf \today}


% Types message, asks the user to type in a command, then
% defines \answer to be the input instead of executing it.
%%% comment out if not needed (next 13 lines)
%\typein[\answer]{Do you want to include comments? (y/n)}
%
%\newcommand{\annotation}[1]
%  {
%  \if\answer y {{\tt #1}}
%  \fi
%  }
%
%\if\answer y
%\typeout{I shall INCLUDE comments.}
%\else
%\typeout{Comments will be NOT shown.}
%\fi

% To see corrections comment next line and uncomment the second one
\newcommand{\correction}[2]{#1}
% \newcommand{\correction}[2]{#2}

\begin{document}

\begin{center}{\bf \Large
\dummy
}
\end{center}

%%%%%%%%%%%%%%%%%%%%%%%%%%%%%%%%%%%%%%%%%%%%%%%%%%%%%%%%%%%%%%%%%%%%%%%%%%%
%                                                                         %
%                      Main part                                          %
%                                                                         %
%%%%%%%%%%%%%%%%%%%%%%%%%%%%%%%%%%%%%%%%%%%%%%%%%%%%%%%%%%%%%%%%%%%%%%%%%%%

A problem with the dominating topology:
Let us consider the following two spaces:
$\omega^\omega$ and $\omega^{\omega\uparrow}$.
Of course, they are homeomorphic by the standard
homeomorphism: 
$\psi \colon \omega^\omega \to \omega^{\omega\uparrow}$ defined
by the following:
$\psi(x)(n) = \sum_{k=1}^n x(k) + n$.

We can consider the dominating topology defined
on the space $\omega^\omega$ (by taking into accout 
the standard basis Hechler neighbourhoods:
$D(N, f) = \{g \in \omega^\omega \colon 
f \restriction N = g \restriction N \wedge
\forall_k f(k) \leq g(k)\}$.) In this way 
we obtain the standard dominating topology
on the space $\omega^\omega$. Using the
standard bijection $\psi$ we can "move" this
topology into the space $\omega^{\omega\uparrow}$.
  However, notice that the space 
$\omega^{\omega\uparrow}$ is a subspace
of the space $\omega^\omega$. So the precise
formulation of my question is like follows:
Does the topologies: The topology on 
the space $\omega^{\omega\uparrow}$
inherited from the usuall dominating topology
defined on $\omega^\omega$ (we treat here the
space $\omega^{\omega\up}$ as a subspace of the space
$\omega^\omega$) and the topology
also defined on the same space $\omega^{\omega\uparrow}$
but ``moved'' from the space $\omega^\omega$ (equipped
with the dominating topology of course) by the 
bijection $\psi$ are the same?
  
  This question is related to the two possibilities 
in introduction of the dominating topology in the space
$[\omega]^\omega$. Indeed, we can ``import'' the dominating topology
into the space $[\omega]^\omega$ by two possible ways:
The first one simply by the homeomorphism between 
$[\omega]^\omega$ and $\omega^\omega$ and the second one, by the
homeomorphism between $[\omega]^\omega$ and
$\omega^{\omega\uparrow}$ which is simply a subset of the space
$\omega^\omega$. In both cases we can move the topology 
by the functions. So we indeed obtain the two topologies.
Simply say that:

  We consider the canonical bijection:
$\Psi \colon [\omega]^\omega \to \omega^\omega$
and by this bijection we can move the dominating topology
from $\omega^\omega$ into the space $[\omega]^\omega$.

the second possibility is to consider the
cannonical bijecton 
$\Xi \colon [\omega]^\omega \to \omega^{\omega\uparrow}$
and then move the topology on $\omega^{\omega\uparrow}$
induced by the domination topology on $\omega^\omega$.
Are these two topologies the same? I think, the answer is
false. However, we can even have that the families of
Baire sets in our topologies are the same. 

OK, so let us formulate these two questions 
in a concise way:

\begin{question}
Let us introduce two topologies on the space
$\omega^\omega$. Define:
$\psi \colon \omega^\omega \to \omega^{\omega\uparrow}$ 
by the following:
$\psi(x)(n) = \sum_{k=1}^n x(k) + n$.
Then we can define the topology:
$\tau_1 = \{ \psi[U]\colon U \in \tau_D\}$, 
where of course $\tau_D$ denotes the dominating topology
on $\omega^\omega$.
However, consider $\omega^{\omega\uparrow}$ as a
subspace of the space $\omega^\omega$ and then 
consider the topology 
$\tau_2 = \{U \cap \omega^{\omega\uparrow}
\colon U \in \tau_D \}$.
The question is: Are these two topologies the same?
Or perhaps they have the same Baire sets?
Or perhaps they have the same meager sets?
\end{question}

\newpage
Take the dominating topology as the base topology on 
$\baire$. \\
\\
Say a set is $\psi$-open if it is the image of 
an open set under $\psi$.\\
\\
{\bf Fact.}
Every relative open set in $\incr$ is $\psi$-open.\\
\bp
Let $U$ be a basic relative open subset of $\incr$, 
say $U = [N, f] \cap \incr$,
where $f$ is in $\baire$.
Note that if $f|N$ is not increasing, then $U$ is empty, 
so we may assume it is increasing. 
%Define
%$$
%F(n) = f(n) \mbox{ if } n\le N, \mbox{ and } F(n) = max\{n,f(n)\} 
%	\mbox{ otherwise.}
%F(0) = f(0) \mbox{ and } F(n+1) = \max\{f(n+1), F(n) + 1 \}.
%$$

%Observe that $U = [N, F]$.
%and $[N,F]$ is an open subset of $\baire$.
%As $\psi$ is a homeomorphism, $U$ is $\psi$-open.

Suppose $g \in U$, and consider $V = [N, \psi^{-1}(g)]$ as a basic open 
subset of $\baire$.
If $h \in V$, then $\forall^\infty n, g(n) \le \psi(h)(n)$, and 
$\forall^\infty n, f(n) \le g(n)$. Therefore $\forall^\infty n,
 f(n) \le \psi(h)(n)$ and hence $\psi(h) \in U$. 
 Thus $g \in \psi (V) \subseteq U$.  It follows that
$$
U = \cup_{g\in U} \;\psi\left([N, \psi^{-1}(g)]\right)
	= \psi\left(  \cup_{g\in U} [N, \psi^{-1}(g)] \right).
$$
Therefore $U$ is $\psi$-open.
\ep

\noindent
{\bf Fact.}
There is a basic $\psi$-open set that's not a relative open set.\\
\bp
Define $f \in\baire$ by $f(2k) = 4k$ and $f(2k+1) = 4k + 3$.
%$g(2k) = 4k + 1$, and $g(2k+1) = 4k + 2$.

Suppose the basic $\psi$-open set 
$$
\psi [N,f] = \bigcup_{h\in \cI} [N_h,h] \cap \incr,
$$
i.e. is relatively open, where $\cI\subseteq \baire$ is a 
suitably chosen index set.
Select $h\in \cI$ such that $\psi(f) \in [N_h,h]$.

Construct $g\in\baire$ by 
$g(2k) = f(2k) + 1$ and $g(2k+1) = f(2k+1) - 1$ whenever
$2k\ge N_h$, and $g(n) = f(n)$ otherwise.
It's not hard to see that $\psi(f) \le \psi(g)$ 
but that neither $f\le g$ nor $g\le f$
on any final segment of $\oo$. (Note: such a $g$ can
be constructed for any $f$ that satisfies $f(n+1) - f(n) \ge 3$
for infinitely many $n$.)


Thus $h \le \psi(f) \le \psi(g)$ implies that 
$\psi(g) \in [N_h, h] \cap \incr \subseteq \psi[N,f]$.
But $g\notin [N,f]$ and $\psi$ is 1-1, a contradiction.
\ep

\noindent
Conclusion:
If these two proofs are correct, then 
the $\psi$-generated topoology
strictly refines the relative topology.\\
\\
Question: $\psi[\emptyset,\overline{0}] = \incr$. 
Is $\psi[N,f]$ dense in $[N,\psi(f)]$ for every $f$?
For what $f$ does $\psi[N,f]$ have {\em any} 
``largeness/smallness'' properties?

\newpage
\newcommand{\inversepsi}{\psi^{-1}}
\section{more stuff}
A {\bf  recursive formula} for $\inversepsi:\incr\to\baire$. 
Supppose that $h\in\incr$. Then
$\inversepsi(h)(0) = h(0)$, and for all $n\ge 1$,
$$
\inversepsi(h)(n) = h(n) - h(n-1) - 1.
$$
\bp
For $n=0$, it's easy to see that $\inversepsi(h)(0) = h(0)$, 
because 
$$
h(0) = \psi(\inversepsi(h)(0))) = \inversepsi(h)(0) + 0.
$$
For $n=1$, $\inversepsi(h)(1) = h(1) - h(0)  - 1$, 
because 
$$
h(1) = \psi(\inversepsi(h))(1) = \inversepsi(h)(0) + \inversepsi(h)(1) + 1
	= h(0) + \inversepsi(h)(1) + 1.
$$
In general suppose it's true for $n<\oo$.
\begin{eqnarray*}
h(n+1) &= &\psi(\inversepsi(h)(n+1)) = \inversepsi(h)(n+1)
	+ \left(\sum_{k=0}^n \inversepsi(h)(k) \right) + n+1 \\
	&=& \inversepsi(h)(n+1) 
	+ \left( h(0) + \sum_{k=1}^n [h(k) - h(k-1) - 1]\right)
	+ n + 1 \\
	&=& \inversepsi(h)(n+1) )+ h(n) + 1.
\end{eqnarray*}
It follows that
$$
\inversepsi(h)(n+1) = h(n+1) - h(n) - 1.
$$
\ep

\noindent
{\bf Another formula.} For $n\ge 0$, \\
$$
\inversepsi(h)(n) = h(n) - 
	\sum_{k=0}^{n-1} \inversepsi(h)(k) - n.
$$
To see this,
%\begin{eqnarray*}
$$
h(n) = \psi(\inversepsi(h))(n) = 
\inversepsi(h)(n) + \sum_{k=0}^{n-1} \inversepsi(h)(k) + n.
%= [h(n) - \sum_{k=0}^{n-1} \inversepsi(h)(k) + n ]
%+ \sum_{k=0}^{n-1} \inversepsi(h)(k) + n.
$$
%\end{eqnarray*}
%As $\psi$ is 1-1, the formula is verified.

\newpage
\noindent
{\bf   Question.}
Suppose that $f\in\baire$ and $h\in\incr$.
Under what conditions is $\inversepsi(h) \ge f$?\\
\\
%For convenience, define $h(-1) = -1$.
%The  
%statement
%$\inversepsi(h)(n) \ge f(n)$
%is equivalent to $h(n) - h(n-1) - 1 \ge f(n)$.
Answer: If 
\begin{equation}
\label{main} 
h(n) - h(n-1) - 1 < f(n) ,
\end{equation}
then $\inversepsi(h)(n) < f(n)$,
and if (\ref{main}) holds for 
infinitely many $n$ then $\inversepsi(h) \not\ge f$
on any final segment of $\oo$.

An earlier example was constructed to show that $\psi([N,f])$
is  not necessarily open. The next proposition generalizes this.\\
\\
{\bf Proposition.} [editted on Oct 16].
For every $f\in \baire$ that doesn't terminate in $0$'s,
the image $\psi([\emptyset,f])$ doesn't
contain an open set.
\bp
Let $A= \{n:\;f(n) \not=0\}$ with the standard increasing enumeration
$(a_k:\;k<\oo)$.
Set $A_{Od} = \{a_{2k+1}:\;k<\oo\}$ and
suppose that $M<\oo$ and $g\in\incr$.
%suppose that $[M,g]$ is a basic open set.
%et $[M,g]\subseteq \psi([\emptyset,f])$ 
%be an arbitrary open set.
%For our purposes we may assume WLOG that $M$ is odd.
Define $h\in [M,g]$ by 
$$
h(n) = \left\{ 
	\begin{tabular}{lll}
	$g(n)$ && if $n < M$\\
	$g(n) + 1$  && if $n \in A_{Od} \setminus M$\\
	$g(n+1)$ && otherwise
	\end{tabular}
	\right.
$$
Suppose $n\in A_{Od} \setminus M$. Then 
$h(n) = g(n) + 1$. And if $n \notin A_{Od}\cup M$,
then $h(n) = g(n+1)$.
In either case, $h(n) \ge g(n)$
and thus $h\in [M,g]$.

For $n\in A_{Od} \setminus (M+1)$,
we have 
$$
\inversepsi(h)(n) = h(n) - h(n-1) - 1 = g(n) + 1 - g(n-1+1) - 1 = 0 < f(n).
$$
%whenever $f(n) \not= 0$. We know that $f(n) \not= 0$
%for infinitely many $n$, and we will assume that 
%$f(n) \not= 0$ for infinitely many odd $n$. 
Thus $(\exists^\infty n)\;\inversepsi(h)(n) < f(n)$,
and it follows that $\inversepsi(h)\notin [\emptyset,f]$.
Therefore $h \in [M,g] \setminus \psi([\emptyset,f])$.
\ep


\section{Added Sat, Oct 16 - Sunday, Oct 17} 
\noindent
%If $f(n) = 0$ for all but finitely many $n\in Od$, then
%we simply switch ``Od'' with ``Ev''  and ``odd'' with ``even''
%throughout, and the proof goes through as before.\\
%\\
Note 1: Since $g$ in the proof could be {\em any} function
in $\incr$, the collection of functions $h
\notin \psi([\emptyset, f])$ is dense.\\
\\
Note 2: The proposition applies to a co-countable subset of $\baire$.\\
\\
{\bf Corollary.} 
If $f\in\baire$ is not eventually zero, then
$\psi([\emptyset, f])$ is co-dense.\\
\\
{\bf Questions.}
Is $\psi([\emptyset,f])$ closed? Is the $\psi$ topology ccc?
Is the $\psi$ topology a Baire space?\\
\\
{\bf Corollary.}
If $\psi([N, f])$ is closed, then it's nowhere dense.
\bp
If $\psi([N, f])$ is closed and dense in some open $[M,g]$,
then it contains $[M,g]$, a contradiction
\ep

\noindent
{\bf Fact.}
The $\psi$ topology is ccc.
\bp
Every $\psi$-open set pulls back to an open set in $\baire$
and the dominating topology is ccc.
\ep

\noindent
{\bf Fact.}
$\psi$ is order preserving but $\inversepsi$ is not.\\
\\
{\bf Fact.} Basic open sets $[N,f]$ and $[M,g]$ are disjoint
if and only if $f|N\perp g|M$.
\bp
The ``if'' direction is trivial, for if 
$f|N\perp g|M$, then $h\in [N,f]$ implies $h|N \not= g|M$.
For the other direction, suppose $f|N$ is compatible with $g|M$ and
WLOG that $f|N \subseteq g|M$ (if not just rename $f$ and $g$).
Define $h|M = g|M$ and $h = f\vee g$ for $n\ge M$. 
Clearly $\forall^\infty n, f(n) \le h(n)$ and $g(n) \le h(n)$.
In addition, $h|N = f|N$ and $h|M = g|M$, so $h \in [N,f]\cap [M,g]$.
\ep

\noindent
{\bf Questions.}
If $f,g\in\baire$ correspond to (almost) disjoint subsets of $\oo$,
what can we say about $\psi(f)$ and $\psi(g)$? What if $f\le g$?\\
\\
Answer: Not much. The particular examples of $f$ and $g$ in a previous proof
correspond to disjoint sets, and $\psi(f) \le \psi(g)$. On the other
hand, it would be easy to construct disjoint $f$ and $g$ such
that $\psi(f)(n) < \psi(g)(n)$ infinitely often and
$\psi(f)(n) > \psi(g)(n)$ infinitely often.
However, if $f\le g$ then $\psi(f)\le \psi(g)$.
What else??\\
\\
{\bf Theorem.}
For all $f\in\baire$ and $N <\oo$, 
$\psi([N, f])$ is dense in $[N, \psi(f)]$,
and if $f$ doesn't terminate in $0$'s, then it's
co-dense as well.
\bp
The part about co-density has been shown above.
To show it's dense,
%For convenience, we assume WLOG that $N\ge 1$.  
take $M \ge N$ and $g\in [N, \psi(f)]$.
We need $h\in[M,g]$ such that 
%$\inversepsi(h)|N = f|N$ and 
$\inversepsi(h) \in [N,f]$.
In other words, $h|M = g|M$ and  
$\forall^\infty n$, 
$$
h(n) \ge g(n) \;\;\mbox{  and  }\;\; h(n) - h(n-1) - 1 \ge f(n).
$$
This is easy to do. Set $h(0) = g(0)$, 
$h|(M+1) = g|(M+1)$ and for $n\ge M+1$, define  
$h(n) = g(n) + h(n-1) + f(n)$.
(We use $n\ge M+1$ to make sure that $n\not=0$.) 
Then for $n\ge M+1$ we have 
$\inversepsi(h)(n) = h(n) - h(n-1) - 1 = f(n) + g(n) - 1 \ge f(n)$.
\ep

\noindent
{\bf Corollary.}
Suppose $f\in \baire$ doesn't terminate in $0$'s and $N<\oo$. Then
$\psi([N,f])$ is neither open nor closed.
\bp
It's not open by a previous result and if 
it was closed then it would contain $[N,\psi(f)]$ since it's dense
in $[N,\psi(f)]$.
This is a contradiction.
\ep

\noindent
{\bf Conjecture.}
$\psi$ is  not continuous, although it is clearly $\psi$-continuous.
\bp
If a basic open set pulls back to an open set, 
then it's $\psi$ image is ?? not open ?? if $f$ isn't eventually $0$.
Clearly it doesn't pull back to a basic open set.

Suppose $\inversepsi([N,g])$ is open. .....
\ep

\noindent
{\bf Conjecture.} The $\psi$-topology is a Baire space.


\end{document}
